\documentclass[10pt,a4paper]{article}
\usepackage[utf8]{inputenc}
\usepackage[T1]{fontenc}
\usepackage[french]{babel}
\usepackage[margin=1.7cm]{geometry}
\usepackage{booktabs}
\usepackage{array}
\usepackage{enumitem}
\usepackage{xcolor}
\setlist{nosep, leftmargin=1.3em}
\setlength{\parindent}{0pt}
\setlength{\parskip}{2pt}
\renewcommand{\arraystretch}{1.15}
\newcommand{\code}[1]{\texttt{\small #1}}
\newcommand{\add}{\textcolor{green!45!black}{\textbf{+}}}
\newcommand{\mod}{\textcolor{orange!75!black}{\textbf{$\sim$}}}
\pagestyle{empty}

\title{\vspace{-1.3cm}\textbf{Intégration de la migration de connexion (TLS) dans OpenFaaS / faasd}\\
\large Fichiers ajoutés et modifiés : rôle et justification}
\author{}\date{}

\begin{document}
\maketitle
\vspace{-1.35cm}

\section{Principe et composants touchés}
Le prototype remplace, sur le chemin \code{/function/<nom>}, le \emph{proxy HTTP} classique du gateway
OpenFaaS par une \textbf{migration de connexion} : le descripteur de la socket TCP du client et
l'\textbf{état de la session TLS~1.3} sont transférés jusqu'au conteneur de la fonction, qui répond
\textbf{directement} au client. Le chemin de données reproduit exactement celui de faasd
(client~$\rightarrow$~gateway~$\rightarrow$~\textbf{faasd-provider}~$\rightarrow$~watchdog~$\rightarrow$~fonction),
mais à chaque saut on \emph{migre} la connexion au lieu de reforwarder du HTTP. Quatre briques sont
concernées : le \textbf{gateway} (\code{faas}), le \textbf{watchdog} (\code{of-watchdog}), le
\textbf{provider} (\code{faasd}), plus deux dépendances de bas niveau : la bibliothèque
\textbf{libtlspeek} et un \textbf{fork de wolfSSL}. Légende : \add~fichier ajouté, \mod~fichier
existant modifié.

\section{Fichiers ajoutés (\add)}

\subsection{Gateway — paquet \code{faas/gateway/httpmigrate/} (entièrement nouveau)}
Ce paquet contient toute la logique de migration côté gateway. \code{main.go} l'active quand
\code{HTTPMIGRATE\_ENABLE=1} (voir §3).

\begin{center}\footnotesize
\begin{tabular}{@{}p{3.15cm} p{13.4cm}@{}}
\toprule
\textbf{Fichier} & \textbf{Rôle / pourquoi il est nécessaire} \\
\midrule
\code{loop.go} & \textbf{Orchestrateur HTTP} (\code{RunLoop}). Ouvre la socket TCP brute, fait un \code{MSG\_PEEK} de la ligne de requête sans la consommer, et route : \code{/function/<nom>}~$\rightarrow$~migration ; sinon~$\rightarrow$~chemin HTTP standard. Nécessaire pour intercepter la connexion \emph{avant} le stack HTTP. \\
\code{loop\_https.go} & \textbf{Orchestrateur HTTPS} (\code{RunLoopHTTPS}), piloté par \textbf{epoll} (une boucle par worker \code{SO\_REUSEPORT}, zéro goroutine par connexion). Termine le TLS, peeke la requête déchiffrée, exporte l'état TLS puis migre. Cœur du mode HTTPS. \\
\code{tls\_listener.\{go,c,h\}} & \textbf{Pont C wolfSSL côté gateway} : machine à états du handshake non bloquant, \code{MSG\_PEEK} déchiffré + \code{wolfSSL\_tls\_export}, et le \code{net.Conn} \code{WolfSSLGtwConn}. Indispensable pour terminer TLS et sérialiser la session ; sert aussi à la terminaison TLS du mode \emph{vanilla} (comparaison équitable). \\
\code{payload.go} & Structure filaire \code{KAPayload} (magic, horodatage \code{top1}, nom de la fonction cible) + (dé)sérialisation. Contrat commun gateway~$\leftrightarrow$~provider~$\leftrightarrow$~watchdog. \\
\code{payload\_https.go} & Charge utile HTTPS (en-tête + état TLS sérialisé) et \code{dispatchMigrateHTTPS} vers le provider. Étend le protocole au cas chiffré. \\
\code{sendfd.go} & \code{sendfd2WithState} : envoi \code{sendmsg}/\code{SCM\_RIGHTS} des FD \code{[clientFD, pipeFD]} + payload en un seul message. Mécanisme noyau de transfert de descripteur. \\
\code{relay\_server.go} & \code{dispatchMigrate} (HTTP) vers le provider ; \code{readCompletionPipe} (lit l'horodatage de fin renvoyé par le watchdog) ; \code{CompletionNotifier} (métriques Prometheus). Relie la migration à l'observabilité du gateway. \\
\code{chan\_listener.go} & \code{ChanListener} : adapte des FD acceptés en \code{net.Listener} pour le \code{http.Server} standard (chemin non-fonction et mode vanilla). \\
\code{conn\_rw.go} & \code{ConnRW} : \code{net.Conn} qui \emph{rejoue} les octets déjà peekés, pour que le serveur HTTP standard relise la requête complète. \\
\code{vanilla\_timing.go} & Middleware et observateur \code{ConnState} de mesure de latence pour le chemin vanilla (instrumentation homogène entre les deux modes). \\
\bottomrule
\end{tabular}
\end{center}

\subsection{Watchdog — \code{of-watchdog/pkg/} (mode « full-proxy »)}
\begin{center}\footnotesize
\begin{tabular}{@{}p{3.15cm} p{13.4cm}@{}}
\toprule
\textbf{Fichier} & \textbf{Rôle / pourquoi il est nécessaire} \\
\midrule
\code{sendfd\_server.go} & Serveur \code{SOCK\_SEQPACKET} du watchdog : reçoit les FD~+~payload migrés et, en mode full-proxy, confie la connexion au pilote. Point d'entrée de la migration côté conteneur. \\
\code{fullproxy.go} & Helpers indépendants du build : transforme chaque requête HTTP en plaintext, la passe au \code{requestHandler} du watchdog (reverse-proxy vers la fonction Python) et re-sérialise la réponse. Sépare \emph{plomberie} et \emph{logique métier}. \\
\code{fullproxy\_cgo.go} & Colle CGO liant \code{wd\_bridge.c} (build \code{wolfssl}) ; expose le callback Go \code{goInvokeHandler}. Nécessaire pour piloter le TLS depuis le C. \\
\code{fullproxy\_stub.go} & Variante sans wolfSSL (build \code{!wolfssl}) : garde la compilation CGO-free du watchdog vanilla. \\
\code{wd\_bridge.\{c,h\}} & \textbf{Pilote de connexion} : restaure la session TLS migrée (\code{tlspeek\_restore}), lit/écrit via wolfSSL, gère le keep-alive, détecte le \emph{wrong-owner} (requête suivante pour une autre fonction) et écrit l'horodatage de fin dans le pipe. Reprend le rôle jadis tenu par les workers C, pour que la fonction ne gère plus aucun FD. \\
\bottomrule
\end{tabular}
\end{center}

\subsection{faasd — \code{faasd/pkg/provider/handlers/migrate\_dispatch.go}}
\textbf{Routeur de migration côté provider} (\code{StartMigrateDispatcher}). Écoute
\code{provider.sock} ; pour chaque message, extrait le nom de fonction cible, résout l'IP CNI du
conteneur \textbf{en interne} (containerd~+~CNI, sans aller-retour HTTP) et retransmet les FD~+~payload
au \code{<ip>.sock} du bon watchdog. Gère aussi le renvoi des connexions \emph{wrong-owner}. Nécessaire
car seul le provider connaît l'IP de chaque conteneur : il joue le même rôle d'aiguillage que dans le
chemin OpenFaaS normal.

\subsection{Dépendances de bas niveau}
\begin{itemize}
\item \textbf{libtlspeek} (\code{lib/tlspeek*.\{c,h\}}, \code{common/}) : bibliothèque C qui \emph{peeke} des
données TLS déchiffrées sans les consommer (\code{tls\_read\_peek}) et qui \textbf{sérialise / restaure}
une session TLS~1.3 entre processus (\code{tlspeek\_serialize}, \code{tlspeek\_restore}). C'est elle qui
rend techniquement possible la migration d'une session chiffrée. (Ses dossiers \code{gateway/},
\code{worker/}, \code{test/} sont la démo et les tests autonomes de la lib, non utilisés par
l'intégration.)
\item \textbf{Fork wolfSSL} (\emph{Migration\_TLS}) : ajoute l'export/import d'état de session
(\code{wolfSSL\_tls\_export} / \code{wolfSSL\_tls\_import}) et l'export des clés de trafic. Sans ces
symboles, une session TLS ne peut pas être reprise dans un autre processus : c'est la pierre angulaire
du prototype.
\end{itemize}

\section{Fichiers existants modifiés (\mod)}
Ces fichiers proviennent d'OpenFaaS / faasd / of-watchdog ; les modifications sont \textbf{additives} et
gardées derrière des variables d'environnement pour ne pas altérer le comportement par défaut.

\begin{center}\footnotesize
\begin{tabular}{@{}p{4.2cm} p{12.35cm}@{}}
\toprule
\textbf{Fichier (existant)} & \textbf{Modification et pourquoi} \\
\midrule
\code{faas/gateway/main.go} &
Point d'entrée : si \code{HTTPMIGRATE\_ENABLE=1}, remplace \code{ListenAndServe()} par
\code{RunLoop}/\code{RunLoopHTTPS} (et par \code{RunVanillaHTTPS} pour la baseline chiffrée), et branche
le \code{CompletionNotifier}. \emph{Pourquoi} : c'est le seul endroit pour dévier le flux d'acceptation
des connexions vers le paquet \code{httpmigrate}, sans forker le serveur HTTP. \\
\code{of-watchdog/config/config.go} &
Ajoute les champs de configuration \code{SendFDEnable}, \code{SendFDSocketDir}, \code{SendFDFullProxy},
\code{OwnFunctionName}, \code{TLSCertFile/KeyFile} et leur lecture depuis l'environnement.
\emph{Pourquoi} : le watchdog doit savoir s'il reçoit des FD, où est le répertoire de sockets partagé,
son propre nom (routage \emph{wrong-owner}) et ses certificats TLS. \\
\code{of-watchdog/pkg/watchdog.go} &
Si \code{SendFDEnable}, démarre \code{StartSendFDServer} en mode full-proxy (ou legacy). \emph{Pourquoi} :
c'est le cycle de vie du watchdog ; il faut y lancer le serveur de réception des FD au démarrage. \\
\code{faasd/cmd/provider.go} &
Si \code{HTTPMIGRATE\_ENABLE=1}, lance \code{StartMigrateDispatcher} (crée \code{provider.sock}).
\emph{Pourquoi} : le provider doit exposer le point d'aiguillage de la migration en même temps que ses
handlers habituels. \\
\code{faasd/.../handlers/deploy.go} &
Bind-mount du répertoire hôte \code{/var/lib/faasd/tlsmigrate} vers \code{/run/tlsmigrate} dans chaque
conteneur. \emph{Pourquoi} : gateway, provider et watchdogs doivent partager les mêmes sockets Unix pour
échanger des FD via \code{SCM\_RIGHTS}. \\
\bottomrule
\end{tabular}
\end{center}

\subsection*{Correctifs de robustesse sur le chemin migré}
Deux corrections ont été nécessaires pour que la session migrée reste correcte sous charge, dans les
fichiers ajoutés ci-dessus : \textbf{(i)} découpe explicite des écritures en records TLS de 16\,384~o
(max app-data TLS~1.3) côté gateway (\code{tls\_listener.go}) et côté watchdog (\code{wd\_bridge.c}), sans
quoi une réponse $>$16~KB provoquait une erreur \emph{record overflow} et une troncature ; \textbf{(ii)}
fermeture explicite du FD brut détaché dans \code{WolfSSLGtwConn.Close} (fuite d'un descripteur par
requête) et \code{wolfSSL\_no\_ticket\_TLSv13} après restauration (éviter d'entrelacer un
\emph{NewSessionTicket} avant la réponse).

\end{document}
